% Appendix: read-out layer of the selection head (VALIDATION; selection procedure, not a result claim). Generated from
% results/metrics/ta_pool_soft_fam.json (eval_curve / by_layer / selector_success; fixed PACO val subset, 500 images,
% 2,288 queries). Needs macros.tex.
\begin{table}[t]
\centering
\setlength{\tabcolsep}{8pt}
\begin{tabular}{l ccc}
\toprule
 & \multicolumn{3}{c}{Read-out layer} \\
\cmidrule(lr){2-4}
Training step & 12 & 24 & 36 \\
\midrule
1k & 52.9 & 76.0 & 75.8 \\
2k & 67.3 & 81.6 & 81.5 \\
3k & 73.4 & 81.9 & 81.9 \\
4k & 75.8 & 82.4 & 82.5 \\
5k & 77.5 & 83.9 & 83.4 \\
6k & 75.8 & 83.1 & 82.7 \\
7k & 77.6 & 83.4 & 83.4 \\
8k & 80.5 & 85.0 & 84.8 \\
9k & 80.5 & 85.5 & 85.6 \\
10k & 80.8 & 85.3 & 85.3 \\
\rowcolor{oursrow}
11k & 81.3 & 85.3 & \textbf{85.6} \\
12k & 81.1 & 85.4 & 85.3 \\
\bottomrule
\end{tabular}
\caption{\textbf{Choice of the read-out layer} (PACO \emph{validation} subset, 500 images; selection success, \%: the
selected mask is a correct one). Selection heads at layers 12, 24 and 36 of the 36-layer decoder are trained jointly;
checkpoint and layer are chosen together on validation (shaded row, bold). Layer 12 is more than 4 points behind at every checkpoint,
while layers 24 and 36 stay within 0.5 points of each other, so the choice among the later layers is not critical.
These numbers are used for model selection only; all results in the paper are on test splits.}
\label{tab:readout_layer}
\end{table}
